\section*{NeurIPS Paper Checklist}

\begin{enumerate}

\item {\bf Claims}
    \item[] Question: Do the main claims made in the abstract and introduction accurately reflect the paper's contributions and scope?
    \item[] Answer: \answerYes{}
    \item[] Justification: The abstract and introduction claim (1) cross-batch metacell learning via affinity-guided prototypes, (2) community structure preservation measured by modularity, (3) improved representation of rare cell types, and (4) capture of spatial microenvironmental structure. All four claims are directly evaluated in the Experiments section with quantitative results.

\item {\bf Limitations}
    \item[] Question: Does the paper discuss the limitations of the work performed by the authors?
    \item[] Answer: \answerYes{}
    \item[] Justification: The Discussion section identifies one limitation: the number of metacells $K$ must be specified in advance, which may not be optimal across datasets and depends on factors such as data quality and scale.

\item {\bf Theory assumptions and proofs}
    \item[] Question: For each theoretical result, does the paper provide the full set of assumptions and a complete (and correct) proof?
    \item[] Answer: \answerNA{}
    \item[] Justification: The paper does not include theoretical results or formal proofs. It proposes an empirical framework and evaluates it experimentally.

\item {\bf Experimental result reproducibility}
    \item[] Question: Does the paper fully disclose all the information needed to reproduce the main experimental results of the paper to the extent that it affects the main claims and/or conclusions of the paper (regardless of whether the code and data are provided or not)?
    \item[] Answer: \answerYes{}
    \item[] Justification: The Method section describes the model architecture, objective functions, and two-stage training procedure in full detail. The Experiments section specifies datasets, affinity construction, baselines, and evaluation metrics. Appendix~\ref{app:training} provides all hyperparameter values (optimizer, learning rate, batch size, loss weights, training duration) and Appendix~\ref{app:calibration} details temperature calibration. All datasets used are publicly available.

\item {\bf Open access to data and code}
    \item[] Question: Does the paper provide open access to the data and code, with sufficient instructions to faithfully reproduce the main experimental results, as described in supplemental material?
    \item[] Answer: \answerNo{}
    \item[] Justification: Code will be made publicly available upon acceptance. All datasets used (Pancreas, PBMC/Immune, Lung, NSCLC) are publicly available and cited in the paper.

\item {\bf Experimental setting/details}
    \item[] Question: Does the paper specify all the training and test details (e.g., data splits, hyperparameters, how they were chosen, type of optimizer) necessary to understand the results?
    \item[] Answer: \answerYes{}
    \item[] Justification: Appendix~\ref{app:training} reports all training and hyperparameter details: optimizer (Adam, lr=$3\times10^{-4}$), batch size (1024), loss weights ($\lambda=0.01$, $\lambda_u=0.1$), training stages (25 epochs pretraining, early stopping for Stage 2), and graph construction settings.

\item {\bf Experiment statistical significance}
    \item[] Question: Does the paper report error bars suitably and correctly defined or other appropriate information about the statistical significance of the experiments?
    \item[] Answer: \answerYes{}
    \item[] Justification: All quantitative results in Tables~1 and~2 are reported as mean~$\pm$~standard deviation across metacells. The source of variability (distribution over metacells) is stated in the table captions.

\item {\bf Experiments compute resources}
    \item[] Question: For each experiment, does the paper provide sufficient information on the computer resources (type of compute workers, memory, time of execution) needed to reproduce the experiments?
    \item[] Answer: \answerNo{}
    \item[] Justification: Compute resource details (GPU type, memory, and approximate training time per experiment) are not currently reported and should be added to the Appendix before submission.

\item {\bf Code of ethics}
    \item[] Question: Does the research conducted in the paper conform, in every respect, with the NeurIPS Code of Ethics \url{https://neurips.cc/public/EthicsGuidelines}?
    \item[] Answer: \answerYes{}
    \item[] Justification: The research uses publicly available, anonymized single-cell transcriptomics datasets for computational method development and raises no ethical concerns.

\item {\bf Broader impacts}
    \item[] Question: Does the paper discuss both potential positive societal impacts and negative societal impacts of the work performed?
    \item[] Answer: \answerNA{}
    \item[] Justification: The paper develops a computational method for single-cell transcriptomics analysis. The positive impact is enabling better biological discovery from multi-batch single-cell datasets. There are no foreseeable negative societal impacts, as the method is not applicable to surveillance, discrimination, or other harmful uses.

\item {\bf Safeguards}
    \item[] Question: Does the paper describe safeguards that have been put in place for responsible release of data or models that have a high risk for misuse (e.g., pre-trained language models, image generators, or scraped datasets)?
    \item[] Answer: \answerNA{}
    \item[] Justification: The paper introduces a method for single-cell biological data analysis. It does not release generative models, scraped datasets, or any asset with significant misuse risk.

\item {\bf Licenses for existing assets}
    \item[] Question: Are the creators or original owners of assets (e.g., code, data, models), used in the paper, properly credited and are the license and terms of use explicitly mentioned and properly respected?
    \item[] Answer: \answerYes{}
    \item[] Justification: All datasets and baseline methods are cited in the paper. The datasets used (Pancreas, PBMC/Immune, Lung, NSCLC) are publicly available research datasets; their original papers are cited. Explicit license names for each dataset should be added to the Appendix before final submission.

\item {\bf New assets}
    \item[] Question: Are new assets introduced in the paper well documented and is the documentation provided alongside the assets?
    \item[] Answer: \answerNo{}
    \item[] Justification: The SCProto codebase will be released upon acceptance with documentation. It is not included in the current submission to preserve anonymity.

\item {\bf Crowdsourcing and research with human subjects}
    \item[] Question: For crowdsourcing experiments and research with human subjects, does the paper include the full text of instructions given to participants and screenshots, if applicable, as well as details about compensation (if any)?
    \item[] Answer: \answerNA{}
    \item[] Justification: The paper does not involve crowdsourcing or research with human subjects. It uses existing, publicly available single-cell transcriptomics datasets.

\item {\bf Institutional review board (IRB) approvals or equivalent for research with human subjects}
    \item[] Question: Does the paper describe potential risks incurred by study participants, whether such risks were disclosed to the subjects, and whether Institutional Review Board (IRB) approvals (or an equivalent approval/review based on the requirements of your country or institution) were obtained?
    \item[] Answer: \answerNA{}
    \item[] Justification: The paper does not involve research with human subjects. All datasets are pre-existing, publicly available single-cell transcriptomics resources.

\item {\bf Declaration of LLM usage}
    \item[] Question: Does the paper describe the usage of LLMs if it is an important, original, or non-standard component of the core methods in this research? Note that if the LLM is used only for writing, editing, or formatting purposes and does \emph{not} impact the core methodology, scientific rigor, or originality of the research, declaration is not required.
    \item[] Answer: \answerNA{}
    \item[] Justification: LLMs are not used as a component of the core SCProto methodology. The framework uses a CVAE with prototype-based representation learning and affinity-guided objectives.

\end{enumerate}
